
\documentclass[10pt]{article}
\usepackage[utf8]{inputenc}
\usepackage[T1]{fontenc}
\usepackage[spanish]{babel}
\usepackage{amsmath}
\usepackage{graphicx}
\usepackage{booktabs}
\usepackage{geometry}
\usepackage{microtype}
\geometry{margin=1.35cm}
\setlength{\parindent}{0pt}
\setlength{\parskip}{3pt}

\begin{document}
\pagestyle{empty}

\begin{center}
{\Large \textbf{Taller 2 - Decaimiento radiactivo con Python}}\\
Cristian Andrés Cárdenas Muñoz \hfill Física Computacional 106018C
\end{center}

\textbf{Contexto físico.}
El decaimiento radiactivo describe la disminución del número de núcleos inestables de una muestra. Para una muestra inicial de
$N_0=10\,000$ núcleos y vida media $t_{1/2}=5$ años, el modelo es
\[
N(t)=N_0e^{-\lambda t}, \qquad
\lambda=\frac{\ln 2}{t_{1/2}}=0.13862944\ \text{año}^{-1}.
\]
La vida media permite comprobar el resultado, porque cada 5 años debe permanecer la mitad de la cantidad anterior.

\textbf{Método.}
Se creó un arreglo de tiempos entre 0 y 30 años con \texttt{numpy.linspace}. La expresión de $N(t)$ se evaluó de forma vectorizada con NumPy, sin usar un ciclo \texttt{for}. Para hallar el tiempo correspondiente a una fracción $f=N/N_0$ se despejó
\[
t=-\frac{\ln(f)}{\lambda}.
\]

\begin{center}
\begin{minipage}{0.47\textwidth}
\centering
\begin{tabular}{rrrr}
\toprule
$t$ (años) & $N(t)$ & referencia & estado\\
\midrule
0  & 10000 & 10000 & PASS\\
5  & 5000  & 5000  & PASS\\
10 & 2500  & 2500  & PASS\\
15 & 1250  & 1250  & consistente\\
\bottomrule
\end{tabular}

\vspace{4pt}
\begin{tabular}{lr}
\toprule
Fracción restante & Tiempo\\
\midrule
50\% & 5.0000 años\\
25\% & 10.0000 años\\
10\% & 16.6096 años\\
\bottomrule
\end{tabular}
\end{minipage}
\hfill
\begin{minipage}{0.49\textwidth}
\centering
\includegraphics[width=\linewidth]{decaimiento_taller2.png}
\end{minipage}
\end{center}

\textbf{Resultado y comparación.}
Los valores calculados en una y dos vidas medias coinciden con las fracciones teóricas de 50\% y 25\%; la celda final del notebook imprime \texttt{PASS} para $t=0$, $5$ y $10$ años. A tres vidas medias permanecen 1250 núcleos, es decir, 12.5\% de la muestra inicial. El 10\% se alcanza aproximadamente a los 16.61 años.

\textbf{Decisión de implementación.}
La principal decisión fue usar operaciones vectorizadas de NumPy para calcular todo el arreglo de $N(t)$ en una sola expresión. Esto evita un ciclo innecesario y hace que el código represente directamente la ecuación física. Para el 10\%, en lugar de buscar el punto visualmente en la gráfica, se despejó el tiempo de la ecuación exponencial.

\textbf{Conclusión.}
El modelo numérico reproduce exactamente el comportamiento esperado de la vida media y permite visualizar la evolución continua de la muestra. Como mejora, el mismo procedimiento puede extenderse a datos experimentales para estimar $\lambda$ y comparar el modelo ideal con mediciones reales.

\textbf{Uso ético de IA.}
Se utilizó IA como apoyo para revisar la organización del código y la redacción, los cálculos y verificaciones se ejecutaron y revisaron antes de la entrega.

\end{document}
